\documentclass[12pt]{article}
\usepackage{lineno}
%\usepackage{ametsoc}
\usepackage{natbib}
\usepackage{a4wide}
\usepackage{longtable}
\usepackage{amsfonts}
\usepackage{amssymb}

\usepackage{color}
\usepackage[pdftex]{graphicx}
\usepackage[pdftex,colorlinks, citecolor=black,linkcolor=black,urlcolor=black]{hyperref}
\definecolor{darkgreen}{rgb}{0.1,0.4,0.1}
\definecolor{darkblue}{rgb}{0.1,0.1,0.3}



\newif\ifdetail
\global\detailfalse
%\global\detailtrue

\def\bem#1{\colorbox{green}{\em #1}}

\newcommand{\beq}  { \begin{eqnarray} }
\newcommand{\eeq}  { \end{eqnarray}}
\newcommand{\beeq}  { \begin{eqnarray*} }
\newcommand{\eeeq}  { \end{eqnarray*}}
\newcommand{\eq }  [1]{Eq.~(\ref{#1})}  % Eq. (#1)
\newcommand{\eqs }  [1]{Eqs.~(\ref{#1})}
\newcommand{\sect}  [1]{Sec.~(\ref{#1})}  % Eq. (#1)
\newcommand{\fig}  [1]{Fig.~\ref{#1}}   % Fig. #1
\renewcommand{\v}[1]{{\mbox{\boldmath$ #1 $}}}  % Vektor durch Fettdruck
\newcommand{\m}    [1]{{\bf  #1 }}              % Matrix in sans serif
\newcommand{\gz}[1] {{{d #1} \over {dz}}}
\newcommand{\pt}[1]{{\partial_t #1}}
\newcommand{\ptt}[1]{{\partial_{tt} #1}}
\newcommand{\pttt}[1]{{\partial_{ttt} #1}}
\newcommand{\pttx}[1]{{{\partial^3 #1} \over {\partial t^2 \partial x}}}
\newcommand{\ptty}[1]{{{\partial^3 #1} \over {\partial t^2 \partial y}}}
\newcommand{\ptx}[1]{{{\partial^2 #1} \over {\partial t \partial x}}}
\newcommand{\pty}[1]{{{\partial^2 #1} \over {\partial t \partial y}}}
\newcommand{\pz}[1]{{\partial_z #1}}
\newcommand{\pzz}[1]{{\partial_{zz} #1}}
\newcommand{\px}[1]{{\partial_x #1}}
\newcommand{\py}[1]{{\partial_y #1}}
\newcommand{\pxi}[1]{{\partial_i #1}}
\newcommand{\pxj}[1]{{\partial_j #1}}
\newcommand{\pxk}[1]{{\partial_k #1}}
\newcommand{\pb}[1]{{{\partial #1} \over {\partial b}}}
%\newcommand{\pxy}[1]{{{\partial^2 #1} \over {\partial x \partial y}}}
\newcommand{\pxy}[1]{\partial_{xy} #1}
\newcommand{\pyy}[1]{\partial_{yy} #1}
\newcommand{\pxx}[1]{\partial_{xx} #1}
\newcommand{\kx}{{\bf k} \times}
\newcommand{\dt}[1]{{{d #1} \over {d t}}}
\newcommand{\Dt}[1]{{{D #1} \over {D t}}}
\newcommand{\Div}[1]{ \nabla \cdot #1 }
\renewcommand{\div}[1]{ \nabla_{yz} \cdot #1 }
\newcommand{\Bar}[1]{ \overline{ #1} }
\newcommand{\rvec} [1]{
    \raisebox{-1.5ex}{$\stackrel{\textstyle #1}{\neg}$} }% rotierter Vektor
\newcommand{\nablq}   {\rvec\nabla}                    % rotated nabla operator
\newcommand{\Vec}[1]{\left(\begin{array}{c} #1 \end{array}\right) }
%\renewcommand{\mho}{\mathcal{W}}
\newcommand{\pn}[1]{{{\partial #1} \over {\partial n}}}
\newcommand{\pnn}[1]{{{\partial #1} \over {\partial m}}}
\newcommand{\ps}[1]{{{\partial #1} \over {\partial s}}}
\newcommand{\p}{{\partial}}
\newcommand{\vn}{{\v \nabla}}
\newcommand{\lk}{\left(}
\newcommand{\rk}{\right)}

\newcommand{\tl}{\tilde}

\newcommand{\I}{\mathcal{I}}
\newcommand{\E}{\mathcal{E}}
\newcommand{\A}{\mathcal{A}}

\newcommand{\ep}{\epsilon}
\newcommand{\erfc}{\mathrm{erfc}}
 

\begin{document}

\section*{The exact scattering integral $S_{nl}$ in wim\_simple: \\ formulation, discretisation and validation}

This note documents the exact evaluation of the non-linear transfer $S_{nl}$
implemented in the sub-package snl\_exact of wim\_simple
(''snonlin\_backend''$\,=\,$'exact').
It complements the model documentation, which contains only a summary,
with the underlying conventions, the derivation chain of the coupling
coefficient, the geometry of the resonance loci, the discretisation on the
spectral grid, the discrete conservation properties, and the validation
against independent implementations.
Only deep water is considered.

\subsection*{Conventions and the kinetic equation}

The kinetic equation below is written for the wave action spectral density
$A(\v k)$ per $d^2 k$. It is related to the surface elevation spectrum
$F(\v k)$, with $\langle \zeta^2 \rangle = \int F \, d\v k$, by
\beq \label{eq:convert}
 A \;=\; \frac{g}{\sigma}\, F   ~~,~~ \sigma(\v k) = \sqrt{g |\v k|}
\eeq
The model carries its spectra as functions of cyclic frequency $\nu$ and
direction $\theta$. Its spectral variables are the variance density
$E(\nu,\theta)$ per Hz and radian and the action $E/\sigma$, which is also a
density per Hz and radian. This model action is not the density $A$ per
$d^2 k$ of \eq{eq:convert}. The model hands $E$ to the exact backend.
With $k = (2\pi\nu)^2/g$ and $dk/d\nu = 8\pi^2 \nu / g$ the backend converts
$E$ to the action density per $d^2 k$ by
\beq
 A(k,\theta) \;=\; \frac{g \, E(\nu,\theta)}{\sigma \, k \, dk/d\nu}
\eeq
and converts the computed transfer back to $E$ with the same factor.

The evolution of $A$ by resonant four-wave interactions is the kinetic
equation of Hasselmann (1962), written here in the Zakharov form
\beq \label{eq:kinetic}
 S_{nl}(\v k_1) = 4\pi \int T^2_{1234}\,
 \delta(\v k_1 + \v k_2 - \v k_3 - \v k_4)\,
 \delta(\sigma_1 + \sigma_2 - \sigma_3 - \sigma_4)\; \mathcal{B}\;
 d\v k_2 \, d\v k_3 \, d\v k_4
\eeq
with the density bracket
\beq \label{eq:bracket}
 \mathcal{B} \;=\; A_3 A_4 (A_1 + A_2) - A_1 A_2 (A_3 + A_4)
\eeq
Two properties of $\mathcal{B}$ are used repeatedly below: it vanishes
identically for constant $A$, and it vanishes on the resonance manifold for
the Rayleigh--Jeans equilibrium $A \propto 1/\sigma$, since
\beq \label{eq:rj}
 \mathcal{B}\big|_{A = c/\sigma} \;=\;
 c^3\, \frac{(\sigma_1+\sigma_2) - (\sigma_3+\sigma_4)}
           {\sigma_1 \sigma_2 \sigma_3 \sigma_4} \;=\; 0
\eeq
Both properties are preserved by the discretisation (see below) and provide
sharp correctness tests.

The kernel $T_{1234}$ in \eq{eq:kinetic} is the canonical four-wave kernel of
the Zakharov equation
\beq \label{eq:zakharov}
 i\, \p_t b_1 \;=\; \int T_{1234}\, b_2^* b_3 b_4\,
 \delta(\v k_1 + \v k_2 - \v k_3 - \v k_4)\, d\v k_2\, d\v k_3\, d\v k_4
\eeq
in the Fourier convention without $2\pi$ prefactors, for which
$\langle b(\v k) b^*(\v k') \rangle = A(\v k)\,\delta(\v k - \v k')$ holds
with the physical action density \eq{eq:convert} directly, and
$T(\v k,\v k,\v k,\v k) = k^3$.
In the convention of Krasitskii (1994), with a factor $1/2\pi$ per transform,
the same kernel reads $T^{K} = T/(2\pi)^2$ and
$T^{K}(\v k,\v k,\v k,\v k) = k^3/4\pi^2$; the collision prefactor $4\pi$
in \eq{eq:kinetic} is form-invariant under this change of convention.
The combination
\beq \label{eq:webb}
 G \;=\; 4\pi\, T^2 \;=\; 4\pi (2\pi)^4 \lk T^{K} \rk^2
\eeq
is the deep-water coupling coefficient of Webb (1978).

\subsection*{Derivation and verification of the kernel}

The kernel is not transcribed from the literature but derived from the
water-wave Hamiltonian, truncated at fourth order in wave steepness,
\beq
 H_3 &=& \frac{1}{2} \int \zeta \lk (\vn \psi)^2 - (G_0 \psi)^2 \rk d\v x \\
 H_4 &=& \frac{1}{2} \int \zeta \, (G_0 \psi)
         \lk G_0 (\zeta\, G_0 \psi) + \zeta \, \nabla^2 \psi \rk d\v x
\eeq
with the surface potential $\psi$, the flat-surface Dirichlet--Neumann
operator $G_0$ (symbol $|\v k|$ in deep water), and normal variables defined by
\beq
 \hat\zeta(\v k) = \sqrt{\frac{k}{2\sigma}} \lk a_{\v k} + a^*_{-\v k} \rk
 ~~,~~
 \hat\psi(\v k) = -i \sqrt{\frac{\sigma}{2 k}} \lk a_{\v k} - a^*_{-\v k} \rk
\eeq
(here written in the $2\pi$-free convention).
Both truncated Hamiltonians were verified exactly against the analytic
kinetic energy of a prescribed harmonic flow, with residuals scaling as the
neglected orders $\varepsilon^3$, $\varepsilon^4$, $\varepsilon^5$.
This check was made during the development. It is not part of the tests
shipped with the sub-package (test\_kernel.py, test\_locus.py,
test\_assembly.py and test\_physics.py).
The canonical (normal-form) transformation then removes the non-resonant
cubic terms; on the four-wave resonance manifold the two classes of induced
quartic contributions coincide pairwise, which is exploited by keeping both
with weight $1/4$ so that the $(12)\!\leftrightarrow\!(34)$ symmetry of $T$
is manifest also off the manifold.

The resulting kernel satisfies, at machine precision, the on-shell
Hermiticity and index symmetries, rotation invariance, cubic homogeneity
$T(\lambda \v k_i) = \lambda^3 T(\v k_i)$, and independence of $g$.
It reproduces the two-wave degenerate-quartet interaction coefficients of
Longuet-Higgins/Leblanc (as compiled by Stuhlmeier, 2024) to about
$10^{-6}$.
These coefficients are values of $T$ at degenerate quadruplets such as
$(\v k, \v k, \v k, \v k)$. There the expression for $T$ contains vanishing
difference wavenumbers and cannot be evaluated directly.
The tests therefore evaluate $T$ at nearby quadruplets
$(\v k_i, \v k_j, \v k_i + \v d, \v k_j - \v d)$ with a small offset
$\v d$ of length $10^{-5}$ or $10^{-6}$ times a fixed vector of order one.
The agreement of about $10^{-6}$ is set by this offset, not by the kernel.
The deviation decreases linearly with the offset and is
$1.2 \cdot 10^{-9}$ for $(\v k,\v k,\v k,\v k)$ at an offset factor of
$10^{-8}$.
During the development the kernel also reproduced the Stokes
second-harmonic and mean-flow response and the bichromatic super- and
sub-harmonic amplitudes. These checks are not part of the shipped tests.
Finally, the identity
\beq \label{eq:xcwebb}
  \mathrm{xc\_webb} \;=\; 64 \pi^5 g^2 \lk T^{K}\rk^2 \;=\; g^2 \, G
\eeq
connecting \eq{eq:webb} to the coupling routine of the exact-NL code of
van Vledder (mod\_xnl4v5, as distributed with WAVEWATCH III and SWAN) was
verified to a relative accuracy of $3 \cdot 10^{-13}$ on resonant quadruplets,
so that the absolute normalisation of \eq{eq:kinetic} is anchored to an
independent production code at the level of the coupling coefficient.

\subsection*{WRT reduction and the resonance loci}

Following Webb (1978) and Tracy and Resio (1982), the wavenumber delta
function in \eq{eq:kinetic} eliminates $\v k_2$, and for fixed
$(\v k_1, \v k_3)$ the frequency delta function reduces the remaining
$\v k_2$-plane integral to a line integral along the resonance locus
\beq \label{eq:locus}
 z(\v k_2) \;=\; q + \sigma(\v k_2) - \sigma(\v k_2 + \v P) \;=\; 0
 ~~,~~ \v P = \v k_1 - \v k_3 ~,~ q = \sigma_1 - \sigma_3
\eeq
with $\v k_4 = \v k_2 + \v P$, so that
\beq \label{eq:wrt}
 S_{nl}(\v k_1) \;=\; \int d\v k_3 \oint_{z=0}
 \frac{ds}{|\vn_{\v k_2} z|}\; G \, \mathcal{B}
\eeq
Note that $z(\v k_3) = 0$ identically: the trivial quadruplet
$\v k_2 = \v k_3$, $\v k_4 = \v k_1$ lies on every locus, with bounded $G$;
the vanishing of the integrand there resides entirely in $\mathcal{B}$.

In deep water the loci form a one-parameter family: scaling lengths by
$|\v P|$ and frequencies by $\sqrt{g |\v P|}$, the locus shape depends only on
\beq
 \hat q = \frac{q}{\sqrt{g |\v P|}} \in (0,1)
\eeq
(the labels $1,3$ are chosen such that $q>0$; $\hat q < 1$ holds for every
pair since $(\sqrt{k_1}-\sqrt{k_3})^2 < k_1 - k_3 \le |\v k_1 - \v k_3|$).
In the normalised frame with $\hat{\v P} = (1,0)$ the two intersections of
the locus with the symmetry axis are analytic,
\beq \label{eq:crossings}
 x_+ = \lk \frac{1-\hat q^2}{2 \hat q} \rk^2
 ~~,~~
 x_- = - \lk \frac{\sqrt{2-\hat q^2} - \hat q}{2} \rk^2
\eeq
which provide exact anchors for the tracing.
The code contains two tracers. Both describe the closed curve by its radius
$r(\theta)$ about the interior midpoint of \eq{eq:crossings}. Here $\theta$
is the polar angle about this midpoint, not the wave direction.
Both compute the derivative of the radius along the curve analytically from
the gradient,
\beq
 r'(\theta) = - \frac{\p z / \p \theta}{\p z / \p r}
 ~~,~~ ds = \sqrt{r^2 + r'^2}\, d\theta
\eeq
The operator uses the routine trace\_batch of assembly.py.
It places the $n_{\rm pts}$ points of a locus uniformly in $\theta$, with
spacing $2\pi/n_{\rm pts}$, and finds the radius of each point by bisection
along its ray.
The periodic trapezoid rule in $\theta$ then yields spectrally accurate
quadrature weights $w_s = (ds/d\theta)\,(2\pi/n_{\rm pts}) / |\vn z|$.
The routine trace\_locus of locus.py is used in the locus tests.
It traces the curve densely by radial continuation from ray to ray and
resamples it to equal arc length $\Delta s$, with the weights
$w_s = \Delta s / |\vn z|$.
For $q = 0$ (equal frequencies, $\v k_1 \ne \v k_3$) the locus degenerates to
the perpendicular bisector $\{ \v k_2 : \v k_2 \cdot \v P = -|\v P|^2/2 \}$;
these same-frequency interactions are unbounded in extent but their line
integral converges through the decay of the spectrum, and they are included
with a separate straight-line quadrature.
The dimensional line measure follows from the scale invariance,
\beq \label{eq:measure}
 \frac{ds}{|\vn z|} \;=\; \frac{|\v P|^{3/2}}{\sqrt{g}}\;
 \frac{d\hat s}{|\hat\vn \hat z|}
\eeq
and the coupling scales as $G \propto |\v k|^6$ by the homogeneity of $T$.
Combining the scalings of \eq{eq:wrt}, the transfer obeys the deep-water
self-similarity
\beq \label{eq:selfsim}
 S_{nl}\lk \v k; \lambda A(\cdot/\mu) \rk \;=\;
 \lambda^3 \mu^{19/2}\, S_{nl}\lk \v k/\mu;  A \rk
\eeq
which the discretisation reproduces on the geometric grid for the closed
loci in the interior of the grid (see below).

\subsection*{Discretisation}

The spectral grid of wim\_simple is geometric in frequency and uniform in
direction, hence geometric in $k$ with uniform spacing $h$ in $\ln k$.
Two consequences are exploited.
First, a rotation of the pair $(\v k_1, \v k_3)$ by a multiple of the
direction spacing $\Delta\theta$ rotates its locus by the same angle.
The locus geometry, the coupling values, the quadrature weights and the
interpolation stencils are therefore tabulated for $\theta_1 = 0$ only.
The table covers every pair $(i_1, i_3)$ of frequency indices and every
direction-index difference $j_3 - j_1$.
Each of these configurations is traced once, exactly.
In the stencils the $\ln k$ indices are absolute grid indices, and the
direction indices are offsets relative to $j_1$.
The table is applied to every $\theta_1$ by shifting all direction indices
by $j_1$.
This renders the discrete operator exactly equivariant under rotations of
the spectrum by grid directions.
The size of the table scales as $n_f^2\, n_d\, n_{\rm pts}$, with $n_f$
frequencies, $n_d$ directions and $n_{\rm pts}$ points per locus.
There is no interpolation in $\hat q$; the discrete grid realises finitely
many configurations, and all of them are solved.
Second, $\hat q$ depends only on the ratio $k_3/k_1$ and the direction
difference.
Let $\rho$ be the ratio of neighbouring frequencies, so that $h = 2\ln\rho$.
A shift of $i_1$ and $i_3$ by one ring multiplies all wavenumbers of the
quadruplets by $\rho^2$, and the stencils move by one index in $\ln k$.
In the interior of the grid this renders the closed-locus part of the
operator equivariant under shifts along the geometric $k$-axis
--- the discrete form of \eq{eq:selfsim}.
Near the grid boundaries this equivariance is broken.
There the strict truncation described below discards different quadruplets
for a spectrum and for its shifted copy.
The same-frequency line loci break it as well.
Each line is cut at $|\v k_2| = 1.02\, k_{max}$, with $k_{max}$ the largest
wavenumber of the grid, and not at a fixed multiple of $|\v P|$.
The shipped test of \eq{eq:selfsim} therefore leaves out the line loci and
evaluates only interior rings.

Action densities at the off-grid locus points $\v k_2, \v k_4$ are obtained
by bilinear interpolation in $(\ln k, \theta)$.
Locus points within a small fraction of $|\v P|$ of the trivial point
$\v k_2 = \v k_3$ are removed: there the exact bracket \eq{eq:bracket}
vanishes but its interpolant does not; the computed transfer is insensitive
to the filter width at the percent level, which is verified as part of the
test battery.
Above the discrete frequency range the spectrum is extended by a parametric
$E \propto \nu^{-5}$ tail over a number of additional geometric rings before
evaluation (following WRT practice in WAVEWATCH III), and the transfer is
restricted to the original band afterwards; quadruplets with any member
beyond the extended range are discarded entirely (strict truncation).
Strict truncation is what preserves the equilibria: every retained quadruplet
carries all four action values, so $\mathcal{B} \equiv 0$ termwise for
constant $A$, and \eq{eq:rj} holds up to interpolation error for
$A \propto 1/\sigma$.
This holds for the operator itself, evaluated without the tail.
The backend in wim\_simple appends the tail by default, with eight rings
(pad\_hi$\,=\,$8, see below).
The tail continues an equilibrium spectrum of the band by
$E \propto \nu^{-5}$, so the extended spectrum is no equilibrium.
The backend therefore returns a nonzero transfer for an equilibrium
spectrum of the band.
For constant $A$ on a band of 20 rings this transfer reaches about $0.2$ of
the one-sided (gross) rate, also in the interior of the band.

With cell areas $\Delta_i = k_i\, \Delta k_i \, \Delta\theta$, where
$\Delta k_i = k_i \lk \rho - 1/\rho \rk = 2 k_i \sinh(h/2)$ is the width of
ring $i$, the discrete form of
\eq{eq:wrt} for one quadruplet instance $(i_1 j_1, i_3 j_3, s)$ carries the
measure
\beq
 \mu_q \;=\; \Delta_{3}\; w_s \; G_s
\eeq
and two assemblies of the same integral are provided.
In the \emph{direct} mode the rate $\mu_q \mathcal{B}$ is accumulated at
$\v k_1$ only --- the discrete analogue of \eq{eq:wrt} and of other WRT
implementations.
In the \emph{conservative} mode each instance deposits the total-action
increments
\beq \label{eq:scatter}
 \lk \delta \mathcal{A}_1, \delta \mathcal{A}_2,
     \delta \mathcal{A}_3, \delta \mathcal{A}_4 \rk
 \;=\; \lk +1,+1,-1,-1 \rk \, \frac{\Delta_1 \mu_q \mathcal{B}}{4}
\eeq
distributed over the grid through the same bilinear stencils used for the
gather.
Dividing the deposit at the interpolated legs by the effective cell area
$\bar\Delta$, interpolated in $\ln k$ with the same weights, conserves the
deposited action exactly, since the stencil weights $w_t$ and the
$\theta$-independent areas satisfy $\sum_t w_t \Delta_{c(t)} = \bar\Delta$
identically.
The factor $1/4$ in \eq{eq:scatter} averages the four role-estimators of the
enumeration over ordered pairs; the total action budget of the discrete
transfer, including the action exchanged with the parametric tail, then
closes at machine precision by construction.
The two modes are independent discretisations of \eq{eq:wrt}; their
difference is $O(h)$, decreases under grid refinement, and serves as a
built-in convergence diagnostic.
Energy and momentum are conserved by the continuous quadruplets and hence to
interpolation accuracy by the discrete transfer.
The code does not compute their residuals.
The backend stores action diagnostics of the last evaluated cell in the
dictionary Snonlin.diag.
In the conservative mode these are action\_rate, the sum of $\Delta_i$
times the transfer over the extended grid, and leak, the action of deposits
that fall outside the grid.
In both modes they are band\_action\_rate and tail\_action\_rate, the same
sum over the band and over the tail rings.

For the semi-implicit time stepping a diagonal derivative is accumulated
alongside,
\beq
 \frac{\p S_{nl}(\v k_1)}{\p A_1} \;=\;
 \sum_q \mu_q \lk A_3 A_4 - A_2 (A_3 + A_4) \rk
\eeq
It is the derivative of the direct rate at $\v k_1$ with respect to $A_1$,
with $A_2$, $A_3$ and $A_4$ held fixed.
The code returns it in both modes.
In the conservative mode, the default, it is not the derivative of the
returned transfer.
There $A_1$ enters also through the deposits of the instances in which
$\v k_1$ plays the role 2, 3 or 4, and the role-1 deposit carries the
factor $1/4$.
At the peak cell of a JONSWAP spectrum and at a cell above the peak the
derivative of the conservative transfer, estimated by finite differences,
differs from the sum above by about a factor 2.
In the direct mode the two differ by about 15\% at these cells, since
$A_1$ also enters through the interpolated $A_2$ and $A_4$.
The sum serves as the diagonal factor of the semi-implicit step.
It is invariant under the conversion between $A$ and $E$ on the diagonal
and is used directly in the update
$E \leftarrow E + S\, \Delta t / (1 - \Delta t\, \p S/\p E)$.

\subsection*{The backend in wim\_simple}

The model selects the exact transfer by
''snonlin\_backend''$\,=\,$'exact'.
In each source step the model hands the variance density $E$ of a cell to
the backend.
The backend returns the transfer $S_{nl}$ of $E$ and the diagonal factor
$\p S/\p E$.
The settings of the backend are attributes of the object Snonlin in the
module pyOM3.wim\_simple.snl\_exact.snonlin\_exact.
They must be set before the setup of the model, since that setup builds the
interaction table.
The attributes and their defaults are
\begin{itemize}
\item g$\,=\,$9.81, the gravitational acceleration of the backend.
It is not taken from the attribute grav of the model.
A model with another value of grav needs the same value here.
\item npts$\,=\,$24, the number $n_{\rm pts}$ of points per closed locus.
The 128 points per locus quoted for the locus tests below refer to the
routine trace\_locus, not to the operator.
\item npts\_line$\,=\,$None, which gives $4\, n_{\rm pts}$ points per
same-frequency line locus.
\item pad\_hi$\,=\,$8, the number of tail rings.
\item dfilt$\,=\,$0.025, the radius of the trivial-point filter in units
of $|\v P|$.
\item cache$\,=\,$None.
A file name here stores the table in that file.
A later setup with the same grid and the same settings reloads it.
\item mode$\,=\,$'conservative', the assembly mode.
The other choice is 'direct'.
\item points$\,=\,$None, the list of cells $(i_x, i_y)$ to evaluate.
None means all cells.
\item max\_cells$\,=\,$64, the largest number of cells per call.
With more cells the backend stops with an error.
\item verbose$\,=\,$True, progress output during the table build.
\end{itemize}
Cells not in points get $S_{nl} = 0$ and $\p S/\p E = 0$.
With the numpy backend of pyOM3 the indices $(i_x, i_y)$ count the interior
cells of the local MPI rank, without halo.
The frequency axis of the model must be geometric.
The setup of the backend checks this and stops otherwise.

\subsection*{Validation}

The implementation is validated layer by layer; the quoted accuracies refer
to the test battery shipped with the sub-package (test grids of 28--30
frequencies and 24 directions).

\begin{itemize}
\item \emph{Kernel}: degenerate two-wave interaction coefficient to $4\cdot
10^{-6}$ at the test offsets; on-shell symmetries, rotation, homogeneity and $g$-independence to
$10^{-13}$; identity \eq{eq:xcwebb} with the WAVEWATCH III coupling routine
to $3\cdot 10^{-13}$.
\item \emph{Loci}: residuals $|z|$ at returned points $4 \cdot 10^{-14}$;
crossings \eq{eq:crossings} exact; resonance closure of dimensional
quadruplets $10^{-14}$ for both signs of $q$; line-delta quadrature against a
brute-force two-dimensional integral to $5\cdot 10^{-5}$ (grid limited);
measure scaling \eq{eq:measure} to $4 \cdot 10^{-15}$; quadrature of
trace\_locus converged to $10^{-9}$ with 128 points per locus.
\item \emph{Assembly}: action budget (transfer $+$ leak) closed to
$10^{-15}$; rotation equivariance, isotropy and cubic homogeneity at
machine precision; discrete self-similarity \eq{eq:selfsim} to
$1.2 \cdot 10^{-7}$ for the closed loci on interior rings;
trivial-filter insensitivity $0.6\%$ between filter widths $0.025$ and
$0.05$.
\item \emph{Equilibria}: constant $A$ null at $2\cdot 10^{-19}$ of the
one-sided (gross) rate; Rayleigh--Jeans null at $7\cdot 10^{-5}$ with $O(h^2)$
convergence; near-vanishing transfer for the Zakharov--Filonenko spectrum
$E \propto \nu^{-4}$.
\item \emph{Magnitude and shape}: mean JONSWAP transfer with the classic
$(+,-,+)$ lobe structure.
The maximum of the positive low-frequency lobe is
$0.8$ times that of the WAM DIA.
The DIA was itself tuned against exact computations
(Hasselmann et al., 1985).
This agreement rules out any residual power of $2\pi$ or $g$ in the
normalisation chain.
\item \emph{Independent exact code}: on identical input, and with the
default interaction filtering of that code disabled, the transfer lobes agree
with the exact-NL code of van Vledder (WAVEWATCH III / SWAN) at the level of
a few percent (positive lobe ratio $0.96$, negative lobe ratio $1.02$ in the
reference case), with the independent result bracketed between the direct and
conservative modes, as expected for consistent discretisations of the same
integral.
\end{itemize}

\subsection*{References}

\noindent Hasselmann, K., 1962: On the non-linear energy transfer in a
gravity-wave spectrum. Part 1. \emph{J. Fluid Mech.}, 12, 481--500. \\
Hasselmann, S., K. Hasselmann, J.H. Allender and T.P. Barnett, 1985:
Computations and parameterizations of the nonlinear energy transfer in a
gravity-wave spectrum. Part II. \emph{J. Phys. Oceanogr.}, 15, 1378--1391. \\
Krasitskii, V.P., 1994: On reduced equations in the Hamiltonian theory of
weakly nonlinear surface waves. \emph{J. Fluid Mech.}, 272, 1--20. \\
Stuhlmeier, R., 2024: Deterministic wave forecasting with the Zakharov
equation. arXiv:2401.03539. \\
Tracy, B.A. and D.T. Resio, 1982: Theory and calculation of the nonlinear
energy transfer between sea waves in deep water. WES Report 11, US Army
Corps of Engineers. \\
van Vledder, G.Ph., 2006: The WRT method for the computation of non-linear
four-wave interactions in discrete spectral wave models.
\emph{Coastal Engineering}, 53, 223--242. \\
Webb, D.J., 1978: Non-linear transfers between sea waves.
\emph{Deep-Sea Res.}, 25, 279--298. \\
Zakharov, V.E., 1968: Stability of periodic waves of finite amplitude on the
surface of a deep fluid. \emph{J. Appl. Mech. Tech. Phys.}, 9, 190--194.

\end{document}
